\chapter*{DANH MỤC KÝ HIỆU VÀ CHỮ VIẾT TẮT}
\addcontentsline{toc}{chapter}{DANH MỤC KÝ HIỆU VÀ CHỮ VIẾT TẮT}

\begin{longtable}{|>{\raggedright\arraybackslash}p{3.6cm}|>{\raggedright\arraybackslash}p{10.4cm}|}
  \hline
  \textbf{Chữ viết tắt} & \textbf{Nghĩa đầy đủ / Giải thích} \\
  \hline
  API   & Application Programming Interface (Giao diện lập trình ứng dụng) \\ \hline
  BLEU  & Bilingual Evaluation Understudy --- hệ số đo độ chính xác $n$-gram có hiệu chỉnh \\ \hline
  BERT  & Bidirectional Encoder Representations from Transformers \\ \hline
  BERTScore & Hệ số tương đồng ngữ nghĩa dựa trên nhúng từ ngữ cảnh (sử dụng PhoBERT trong khóa luận này) \\ \hline
  CAPI  & (Server-side) Conversion API \\ \hline
  CoT   & Chain-of-Thought (Chuỗi suy luận từng bước) \\ \hline
  CSV   & Comma-Separated Values (Dữ liệu phân cách bằng dấu phẩy) \\ \hline
  CRUD  & Create, Read, Update, Delete (Tạo, Đọc, Cập nhật, Xóa) \\ \hline
  DOM   & Document Object Model \\ \hline
  GPT   & Generative Pre-trained Transformer \\ \hline
  HF    & Hugging Face \\ \hline
  JSON  & JavaScript Object Notation \\ \hline
  JSDOM & JavaScript DOM implementation in Node.js \\ \hline
  LLM   & Large Language Model (Mô hình ngôn ngữ lớn) \\ \hline
  MoA   & Mixture-of-Agents --- Kiến trúc đa tác nhân phối hợp (Wang et al., 2024) \\ \hline
  MV3   & Manifest V3 (Đặc tả tiện ích mở rộng Chrome) \\ \hline
  NLP   & Natural Language Processing (Xử lý ngôn ngữ tự nhiên) \\ \hline
  PRD   & Product Requirements Document (Tài liệu yêu cầu sản phẩm) \\ \hline
  RAG   & Retrieval-Augmented Generation (Tạo văn bản tăng cường truy xuất) \\ \hline
  RLS   & Row-Level Security --- Bảo mật cấp hàng (PostgreSQL) \\ \hline
  ROUGE & Recall-Oriented Understudy for Gisting Evaluation \\ \hline
  SAG   & Search-Augmented Generation (Tạo văn bản tăng cường tìm kiếm) \\ \hline
  SDK   & Software Development Kit (Bộ công cụ phát triển phần mềm) \\ \hline
  SSE   & Server-Sent Events (Sự kiện đẩy từ máy chủ) \\ \hline
  TS    & TypeScript \\ \hline
  UI    & User Interface (Giao diện người dùng) \\ \hline
  URL   & Uniform Resource Locator (Địa chỉ tài nguyên thống nhất) \\ \hline
  UUID  & Universally Unique Identifier (Mã định danh duy nhất toàn cầu) \\ \hline
  ViT5  & Vietnamese Text-to-Text Transfer Transformer \\ \hline
  PhoBERT & Mô hình BERT tiền huấn luyện cho tiếng Việt (vinai/phobert-base) \\ \hline
  \multicolumn{2}{l}{} \\
  \hline
  \textbf{Thuật ngữ} & \textbf{Định nghĩa} \\ \hline
  Tác nhân đề xuất (Proposer)    & LLM ở tầng thứ nhất trong pipeline MoA, có nhiệm vụ tạo ra một bản thảo tóm tắt ứng viên \\ \hline
  Tác nhân tổng hợp (Aggregator)  & LLM ở tầng thứ hai trong pipeline MoA, có nhiệm vụ dung hợp các bản thảo thành đầu ra cuối cùng \\ \hline
  Kết nối phần dư (Residual connection) & Việc đưa trực tiếp bài viết nguồn vào ngữ cảnh của tác nhân tổng hợp để quá trình dung hợp luôn neo vào nội dung gốc, không bị ảo tưởng \\ \hline
  Đánh giá rubric   & Hình thức LLM-giám khảo chấm điểm tóm tắt trên nhiều tiêu chí có tên cụ thể theo thang số \\ \hline
  Đánh giá từng cặp (Pairwise) & Hình thức LLM-giám khảo so sánh hai bản tóm tắt A và B, chọn ra bản thắng cuộc \\ \hline
  Hệ số Fleiss' $\kappa$ & Chỉ số đồng thuận giữa các người đánh giá dùng cho dữ liệu phân loại từ $\geq 2$ người đánh giá \\ \hline
  Kiểm định dấu (Sign test)        & Kiểm định phi tham số cho kết quả nhị phân theo từng cặp, sử dụng để kiểm tra tỷ lệ thắng từng cặp \\ \hline
  Phân nhóm theo độ dài (Length bucketing) & Phân tầng các phán quyết theo cặp dựa trên tỷ lệ độ dài đầu ra để kiểm soát thiên kiến độ dài \\ \hline
  Kéo theo logic (Entailment)       & Quan hệ logic: giả thuyết được hỗ trợ bởi tiền đề; sử dụng để đánh giá tính xác thực cấp tuyên bố \\ \hline
\end{longtable}
